\newpage
\section{Literature Review}
\subsection{Business Review / Software Review}
\subsubsection*{Calendly}
\begin{center}
    \includegraphics[width=0.7\textwidth]{images/CalendarConnections.png}
\end{center}
\textbf{Calendly}\cite{calendly_ref} is an automated scheduling platform that eliminates back-and-forth emails by allowing users to share their availability links, enabling invitees to instantly book open time slots.
\subsubsection*{Savvycal}
\begin{center}
    \includegraphics[width=0.7\textwidth]{images/Savvycall.png}
\end{center}
\textbf{Savvycal}\cite{savvycal_ref} is a collaborative scheduling tool featuring an interactive calendar overlay, which allows recipients to lay their own calendar over the sender's schedule to easily find mutual availability.
\subsubsection*{Doodle}
\begin{center}
    \includegraphics[width=0.7\textwidth]{images/doodle-img.png}
\end{center}
\textbf{Doodle}\cite{doodle_ref} is a group scheduling platform that utilizes a polling system, where an organizer proposes multiple time options and participants vote on their availability to determine the best consensus time for a meeting.

\subsubsection*{Comparative Table}
\vspace{0.5 cm}
\input{tables/business_comparison}

\newpage
\subsection{Architectural Foundations \& Prior Art}
Standard scheduling applications only look at calendar availability, they cannot process physical travel times or complex overlapping rules. To solve this, Converge is built upon established Artificial Intelligence and Operations Research methodologies specifically designed for complex logistics.

\subsubsection{Constraint Satisfaction Modeling for Complex Timetabling}
Timetabling problems are formally modeled as Constraint Satisfaction Problems (CSPs), where a set of variables must be assigned values from defined domains while satisfying a collection of constraints. In the context of course scheduling, variables typically represent assignments such as (instructor, time slot, branch), while constraints enforce rules such as no instructor or room being allocated to more than one class at the same time. This formulation captures the combinatorial nature of scheduling, which has been widely classified as an NP-hard problem, meaning that the computational complexity grows exponentially with the number of variables and constraints\cite{elsakka2015}.

In addition to hard constraints, real-world systems also include soft constraints such as preferred schedules or workload balancing, which are incorporated into an objective function for optimization. Furthermore, educational timetabling problems vary significantly depending on institutional rules, resource limitations, and operational requirements, making each real-world scenario highly unique and difficult to generalize\cite{ceschia2022}.

The Converge system adopts this formal CSP/COP framework by modeling scheduling as a structured decision problem, where instructor availability, branch capacity, and time slots form the variable domains, and constraints enforce feasibility. The system’s Constraint Satisfaction AI Engine operates by searching for feasible assignments that satisfy all hard constraints while producing a high-quality schedule under operational conditions.
\subsubsection{Travel Buffer System}
Classical timetabling formulations primarily focus on logical constraints such as time conflicts and resource allocation; however, real-world scheduling introduces additional constraints that are not always explicitly modeled in standard formulations. Research highlights that practical timetabling systems must incorporate domain-specific constraints to reflect real operational environments\cite{schaerf1999}. In the case of multi-branch tutoring centers, one critical constraint is the physical movement of instructors between locations.

To address this, the Converge system extends the traditional CSP model by introducing a \textit{Travel Buffer} constraint. This constraint enforces a minimum temporal separation between consecutive assignments of the same instructor across different branches. Formally, if an instructor is assigned to branch $A$ at time $t$, and branch $B$ requires a travel duration $\Delta$, then any subsequent assignment at branch $B$ must satisfy: $t_2 \geq t + \Delta$

Rather than relying on dynamic external data sources, the system uses a predefined travel time matrix between branches. This design choice ensures deterministic behavior, reduces computational overhead, and maintains consistency with operational policies. By integrating travel constraints directly into the scheduling model, the system prevents infeasible schedules that would otherwise satisfy only abstract constraints but fail in real-world execution.

\subsubsection{Heuristic Search and Interactive Explainable AI (XAI)}
Given the NP-hard nature of timetabling problems, exact algorithms are generally infeasible for large-scale instances. As a result, most practical systems rely on heuristic and metaheuristic search strategies to efficiently explore the solution space. Prior research identifies approaches such as greedy heuristics, local search, and constraint-based methods as effective techniques for constructing feasible timetables within reasonable time constraints\cite{schaerf1999}.

A common heuristic principle in timetabling is to prioritize the most constrained variables first, thereby reducing the likelihood of conflicts later in the search process. In the Converge system, heuristic search strategies are used within the Constraint Satisfaction AI Engine to incrementally build feasible schedules while respecting constraints such as availability, capacity, and travel buffers.

In addition, modern timetabling systems recognize the importance of human involvement in the scheduling process, as fully automated solutions may not capture all practical considerations (A Survey of Automated Timetabling). Therefore, Converge incorporates Explainable AI (XAI) principles by providing clear reasoning behind scheduling decisions. Each assignment is accompanied by a traceable explanation of the constraints that influenced it, allowing administrators to understand, validate, and adjust the generated schedules when necessary.

\subsection{Tool and Technology Review}
\subsubsection{Frontend Language}
\noindent
\begin{minipage}{0.15\textwidth}
    \centering
    \includegraphics[width=1\linewidth]{images/TypeScript.png}
\end{minipage}
\hfill
\begin{minipage}{0.82\textwidth}
We chose \textbf{TypeScript}\cite{typescript_ref} over JavaScript as the frontend language because its compile-time error checking prevents structural logic failures in our multi-branch scheduling. In making this choice, we accept the tradeoff of a slower initial prototyping phase due to strict type definition requirements.
\end{minipage}
\vspace{0.5cm}\\
\input{tables/frontend_language}

\newpage
\subsubsection{Frontend Framework}
\noindent
\begin{minipage}{0.15\textwidth}
    \centering
    \includegraphics[width=1\linewidth]{images/Vue.js.png}
\end{minipage}
\hfill
\begin{minipage}{0.82\textwidth}
We chose \textbf{VueJS}\cite{vuejs_ref} over React because its lower learning curve and built-in tooling allow us to focus on building features rather than configuring the project. In making this choice, we accept the tradeoff of a smaller ecosystem.
\end{minipage}
\vspace{0.5cm} \\
\input{tables/frontend_framework}

\subsubsection{Backend Language}
\noindent
\begin{minipage}{0.15\textwidth}
    \centering
    \includegraphics[width=1\linewidth]{images/go.png}
\end{minipage}
\hfill
\begin{minipage}{0.82\textwidth}
We chose \textbf{Golang}\cite{golang_ref} over Java and NodeJS as the backend language because its compiled performance and native concurrency model perfectly balance raw execution speed with a low memory footprint. In making this choice, we accept the tradeoff of Go’s verbose syntax and manual error handling, which requires writing more boilerplate code than Node.js.
\end{minipage}
\vspace{0.5cm} \\
\input{tables/backend_language}

\subsubsection{Backend Framework}
\noindent
\begin{minipage}{0.15\textwidth}
    \centering
    \includegraphics[width=1\linewidth]{images/gin.png}
\end{minipage}
\hfill
\begin{minipage}{0.82\textwidth}
We chose the \textbf{Gin}\cite{gin_ref} framework over Go’s native standard library (\texttt{net/http}) because its built-in routing and middleware capabilities significantly accelerate the development of our REST API. In making this choice, we accept the tradeoffs of framework lock-in and a slight memory overhead.
\end{minipage}
\vspace{0.5cm} \\
\input{tables/backend_framework}

\subsubsection{Database}
\noindent
\begin{minipage}{0.15\textwidth}
    \centering
    \includegraphics[width=1\linewidth]{images/PostgresSQL.png}
\end{minipage}
\hfill
\begin{minipage}{0.82\textwidth}
We chose \textbf{PostgreSQL}\cite{postgresql_ref} as the database because its native support for range data types is perfectly suited for a scheduling system. In making this choice, we accept the tradeoff of a steeper learning curve and slightly more configuration overhead.
\vspace{0.5cm} \\
\end{minipage}
\input{tables/database}

\subsubsection{Architectural Pattern}
\noindent
\begin{minipage}{0.15\textwidth}
    \centering
    \includegraphics[width=1\linewidth]{images/hexagonal.png}
\end{minipage}
\hfill
\begin{minipage}{0.82\textwidth}
We chose \textbf{Hexagonal Architecture}\cite{hexagonal_ref} as the backend architectural pattern because it keeps the core scheduling domain portable and testable in isolation from the web framework. In making this choice, we accept the tradeoff of increased initial complexity, requiring more folders and interfaces upfront.
\end{minipage}
\vspace{0.5cm} \\
\input{tables/architecture}

\subsubsection{Deployment Architecture}
\noindent
\begin{minipage}{0.15\textwidth}
    \centering
    \includegraphics[width=1\linewidth]{images/modular_monolith.png}
\end{minipage}
\hfill
\begin{minipage}{0.82\textwidth}
We chose a \textbf{Modular Monolith}\cite{architecture_ref} over both Microservices and a standard Monolith to provide a safe evolution path for the system. This approach allows us to properly model the system using Domain-Driven Design (DDD) principles without the immediate operational complexity of a distributed network. In making this choice, we accept the tradeoff of relying on strict developer code discipline to enforce module boundaries, rather than relying on physical infrastructure separation.
\end{minipage}
\vspace{0.5cm} \\
\input{tables/deployment}

\newpage
\subsubsection{Containerization}
\noindent
\begin{minipage}{0.15\textwidth}
    \centering
    \includegraphics[width=1\linewidth]{images/Docker.png}
\end{minipage}
\hfill
\begin{minipage}{0.82\textwidth}
We chose \textbf{Docker}\cite{docker_ref} for containerization over traditional bare-metal because it keeps the setup environment consistency across development. By packaging the application and its dependencies into a single image, we eliminate environment-specific bugs. In making this choice, we accept the tradeoff of an increased learning curve, as the entire development team must understand Docker.
\end{minipage}
\vspace{0.5cm} \\
\input{tables/docker}
